%% Chapitre 10 — Connecter des outils externes avec MCP
%% RÈGLE : uniquement des commandes sémantiques bookforge.
%% CHAPITRE DATÉ : l'écosystème MCP évolue vite.

\chapitre{Connecter des outils externes avec MCP}

\introChapitre{
  Un agent gagne en puissance quand il peut parler à tes autres outils : base
  de données, gestionnaire de tickets, documentation interne. Ce chapitre, daté
  car l'écosystème évolue, explique le principe du protocole MCP et comment
  brancher une première intégration utile sans se faire piéger par les questions
  de confiance.
}

\section{Le problème que résout MCP}

\paragraphe{Jusqu'ici, \code{Claude Code} agissait dans ton dépôt : il lisait
des fichiers, lançait des commandes, modifiait du code. Le chapitre précédent a
industrialisé ces gestes internes avec les slash commands et les hooks. Mais le
travail réel déborde du dépôt. Le ticket qui décrit le bug vit dans Jira ou
Linear. Le schéma exact d'une table vit dans ta base de données. La règle
métier qui justifie une validation vit dans une page de documentation interne.}

\paragraphe{Sans accès à ces sources, l'agent devine. Tu lui colles à la main un
extrait de ticket, tu lui recopies une structure de table, tu lui décris une
règle de mémoire. C'est du travail manuel, et chaque copier-coller approximatif
augmente le risque d'hallucination : l'agent comble les trous avec du plausible.}

\paragraphe{\important{MCP résout exactement ce problème : donner à l'agent un
accès direct, structuré et réutilisable à des outils externes.} Au lieu de
décrire le monde à l'agent, tu le branches dessus.}

\section{Le principe : serveurs, outils, et l'agent qui appelle}

\paragraphe{\terme{MCP} signifie \etranger{Model Context Protocol} : un protocole
standard qui décrit comment un agent dialogue avec des sources externes. Standard
est le mot clé. Avant lui, chaque intégration était du sur-mesure ; aujourd'hui
n'importe quel outil qui parle MCP se branche de la même façon.}

\paragraphe{La brique concrète est le \terme{serveur MCP}. C'est un petit
programme qui se place entre l'agent et un système précis (ta base, ton
gestionnaire de tickets) et qui expose des \emphase{outils} : des actions nommées
avec des paramètres clairs, comme \code{rechercher_ticket} ou
\code{executer_requete_sql}. L'agent ne voit pas la plomberie ; il voit une liste
d'outils disponibles, leur description, et il les appelle quand la tâche le
demande.}

\paragraphe{Le cycle est simple. Tu configures un serveur, \code{Claude Code} se
connecte au démarrage de la session et récupère la liste des outils exposés.
Pendant ton travail, si l'agent juge qu'un outil l'aide, il l'appelle, lit le
résultat, et continue son raisonnement avec cette information fraîche dans son
contexte. Tu n'as rien à recopier.}

\paragraphe{Un serveur tourne en local, lancé comme un sous-processus sur ta
machine, ou à distance, exposé derrière une adresse réseau. Le local convient
pour un outil qui agit sur ton poste ; le distant, pour un service partagé par
toute l'équipe. Dans les deux cas, l'agent les voit de la même manière : une
liste d'outils. C'est tout l'intérêt d'un protocole standard, tu raisonnes en
\emphase{capacités}, pas en tuyauterie.}

\section{Brancher un premier serveur MCP}

\paragraphe{La configuration tient dans un fichier ou dans la commande
\code{claude mcp add}. Tu déclares chaque serveur : un nom, son mode de
connexion, et ses éventuels paramètres. Voici un modèle d'intégration de
recherche documentaire distante. Les valeurs en majuscules sont à remplacer par
celles de ton équipe ; ce n'est pas un paquet public à installer tel quel.}

\begin{extraitcode}{json}
{
  "mcpServers": {
    "docs-interne": {
      "type": "http",
      "url": "https://MCP_DOCS_DE_TON_EQUIPE.example/mcp",
      "headersHelper": "./scripts/mcp-docs-headers.sh"
    }
  }
}
\end{extraitcode}

\paragraphe{Le script \code{headersHelper} fournit les en-têtes d'authentification
au moment de l'appel. Le secret ne vit donc pas en clair dans le fichier de
configuration. On verra plus loin pourquoi ce réflexe compte. Une fois le
serveur ajouté, tu vérifies qu'il est bien vu par l'agent.}

\begin{extraitcode}{bash}
# Ajouter un serveur MCP HTTP au niveau du projet
claude mcp add --transport http docs-interne --scope project https://MCP_DOCS_DE_TON_EQUIPE.example/mcp

# Lister les serveurs MCP connus de la session
claude mcp list

# Démarrer une session et inspecter les outils chargés
claude
> /mcp
\end{extraitcode}

\paragraphe{Si le serveur apparaît avec ses outils, c'est branché. Demande alors
quelque chose de concret : \emphase{cherche dans la doc interne la politique de
nommage des branches}. L'agent appelle l'outil de recherche au lieu d'inventer.}

\begin{avertissement}
Les noms exacts de commandes, de fichiers de configuration et de paquets MCP
changent vite : l'écosystème est jeune et bouge à chaque version de
\code{Claude Code}. Considère les exemples ci-dessus comme une illustration du
\emphase{principe}, pas comme une référence figée. Le réflexe durable : lance
\code{claude mcp --help} et consulte la documentation de la version que tu as
réellement installée avant de copier une configuration trouvée ailleurs.
\end{avertissement}

\section{Choisir des intégrations fiables et limiter les permissions}

\paragraphe{Brancher un serveur MCP, c'est inviter du code tiers à agir au cœur
de ta session. Avant d'en ajouter un, pose-toi deux questions : \emphase{qui l'a
écrit}, et \emphase{de quoi a-t-il vraiment besoin}.}

\paragraphe{Pour la provenance, privilégie les serveurs officiels d'un éditeur ou
ceux d'une source largement utilisée et auditée. Un dépôt inconnu avec trois
étoiles, qui demande un jeton d'accès large, mérite la méfiance. Lis ce qu'il
exécute avant de lui confier tes identifiants.}

\paragraphe{Pour les besoins, applique le principe du moindre privilège, qui
prolonge directement les permissions vues au chapitre des garde-fous. Un serveur
de tickets a besoin de lire et commenter, rarement de supprimer. Un serveur de
base de données peut souvent se contenter d'un accès en lecture seule.}

\paragraphe{Cadrer un serveur ne s'arrête pas au compte qu'il utilise : tu peux
aussi décider, côté agent, quels de ses outils tu autorises sans confirmation. Un
outil de lecture peut tourner librement ; un outil qui écrit ou supprime devrait
rester soumis à ton aval explicite, exactement comme une commande shell sensible.
Le serveur définit ce qui est \emphase{possible} ; tes permissions définissent ce
qui passe \emphase{sans te demander}.}

\begin{astuce}
Crée un compte ou un jeton dédié à chaque serveur MCP, avec le strict minimum de
droits, plutôt que de réutiliser ton compte personnel d'administrateur. Si un
serveur dérape ou fuit son jeton, les dégâts restent bornés à ce périmètre, et tu
peux le révoquer sans casser le reste de ton outillage.
\end{astuce}

\section{Cas concrets : tickets, base de données, documentation}

\paragraphe{Le \emphase{gestionnaire de tickets} est souvent la première
intégration rentable. Tu donnes une référence et l'agent va lire le ticket
lui-même : description, commentaires, critères d'acceptation. Demande
« implémente ce que décrit le ticket PROJ-241 », et il récupère le contexte avant
de toucher au code. Limite-le à la lecture et au commentaire ; ne le laisse pas
fermer ou réassigner sans ton aval.}

\paragraphe{La \emphase{base de données} évite à l'agent d'halluciner un schéma.
Au lieu de supposer qu'une colonne s'appelle \code{user_id}, il interroge la
structure réelle et écrit une requête correcte du premier coup. Branche-le en
lecture seule, sur une réplique ou une base de développement, jamais directement
sur la production en écriture.}

\paragraphe{La \emphase{recherche documentaire} ancre l'agent dans tes décisions
maison. Conventions de code, choix d'architecture, runbook d'incident : l'agent
cite la source au lieu d'inventer une bonne pratique générique. C'est un
complément naturel à la mémoire projet, pour ce qui est trop volumineux ou trop
mouvant pour tenir dans un fichier d'instructions.}

\section{Risques de confiance et bonnes pratiques de cadrage}

\paragraphe{Un serveur MCP n'est pas un simple lecteur passif : il peut renvoyer
du texte qui influence l'agent. Un ticket, une page de doc ou une ligne de base
peuvent contenir des instructions glissées par un tiers. L'agent, qui lit tout
comme du contexte, risque de les suivre. C'est le cœur du problème de confiance.}

\begin{avertissement}
Ne branche que des serveurs MCP de confiance, et traite toute donnée externe
comme potentiellement hostile. Un serveur malveillant ou compromis peut exfiltrer
ton code, vider un jeton d'accès ou pousser l'agent à exécuter une commande
dangereuse via une instruction cachée dans un ticket ou une page web. Trois
garde-fous : n'installe que des serveurs dont tu connais l'auteur, donne-leur le
moins de droits possible, et garde la confirmation des actions sensibles plutôt
que de tout autoriser d'un bloc.
\end{avertissement}

\paragraphe{Le bon cadrage n'est pas de tout interdire, mais de doser. Commence
avec une seule intégration, en lecture seule, sur des données non critiques.
Observe comment l'agent s'en sert, vérifie ses appels d'outils comme tu vérifies
ses diffs, puis élargis le périmètre une fois la confiance établie.}

\resumeChapitre{
  \element{\important{MCP} : un protocole standard qui permet à l'agent d'appeler des outils externes (tickets, base de données, documentation) sans que tu recopies le contexte à la main.}
  \element{\important{Serveur MCP} : un petit programme qui expose des actions nommées — l'agent voit une liste de capacités, pas la plomberie derrière.}
  \element{\important{Moindre privilège} : crée un compte dédié par serveur avec le strict minimum de droits — lecture seule là où l'écriture n'est pas nécessaire.}
  \element{\important{Données externes comme contexte hostile} : un ticket ou une page de doc peut contenir des instructions glissées par un tiers — n'installe que des serveurs de confiance.}
  \element{\important{Progressivité} : commence par une seule intégration en lecture sur des données non critiques, observe les appels d'outils comme tu lis un diff, puis élargis.}
}

\paragraphe{Avec MCP, ton agent cesse d'être enfermé dans le dépôt : il touche à
tout ton environnement de travail. Cette portée nouvelle appelle une nouvelle
montée en puissance. Jusqu'ici, un seul agent menait toute la tâche ; au chapitre
suivant, on apprend à en faire travailler plusieurs ensemble, à déléguer des
parties d'un chantier à des subagents et à orchestrer leur effort.}
